% Section 9 — Related work (charter §4.9; issue T4-2, WOV-275)
\section{Related work}
\label{sec:related}

% Citation hygiene (guardrail G5, from R6, VERBATIM):
%   - Do NOT cite "SOAP" as a work.
%   - Do NOT mischaracterize INDICT (Le et al., NeurIPS 2024).
%   - Verify Clover, Plan-and-Solve, PAL before citing.
%   - Find correct references independently rather than reusing prior
%     confabulated exemplars.
% Every \cite in this section was verified against a primary source at
% insertion time (see theory.bib header). The section positions the paper; it
% does not survey. Each area is cited only where it touches a specific result,
% and no individual ingredient is claimed as this paper's novelty (guardrail
% G2); the circumscribed claim is stated in \cref{subsec:rw-novelty}.

This paper sits at the meeting point of four literatures: information-theoretic
accounts of what a representation can carry, learning-theoretic accounts of what
a learner can extract from it, the categorical machinery we borrow to state
parity exactly, and the empirical study of code representations for language
models that motivates the whole question. We situate each in turn, and in each
case the relationship is the same in form: the prior work supplies an object or a
tradeoff, and our contribution is to fix one axis (target information) so the
other (a bounded learner's value) can be isolated.

\subsection{Information-theoretic views of representation}
\label{subsec:rw-info}

The negative half of the paper (\cref{sec:parity}) is built from sufficiency and
the data-processing inequality, both standard \citep{coverthomas2006,
lehmanncasella1998}: a sufficient statistic for $Y$ retains all predictive
information, and post-processing cannot increase mutual information. The
information-bottleneck programme \citep{tishby2000informationbottleneck}, and its
variational deep-learning realisation \citep{alemi2017deepvib}, turn this calculus
into a design objective by trading compression of $X$ against retention of
information about $Y$ along a Lagrangian frontier. Our setting is the
complementary one. We do not move along that frontier; we hold the predictive
term $\MI{Y}{\rep(X)}$ \emph{fixed} across representations (information parity,
\cref{def:parity}) and ask what is left to vary. The irrelevance theorem
(\cref{thm:irrelevance}) is, in bottleneck language, the statement that two points
on the same level set of retained information are indistinguishable to the Bayes
optimum, so any residual effect cannot be an information-bottleneck effect at all.

\subsection{Learning-theoretic analyses of representation}
\label{subsec:rw-learning}

A second body of work studies, within statistical learning theory, when a
\emph{learned} representation provably helps: multitask representation learning
quantifies the sample-complexity benefit of a shared feature map
\citep{maurer2016multitask}, and transfer-learning theory shows the benefit hinges
on task diversity \citep{tripuraneni2021taskdiversity}. These results ask which
representation to \emph{learn} and prove that the right one lowers estimation
error. The positive half of this paper (\cref{sec:bounded}) asks a different
question on the same machinery: given two representations of equal information,
what advantage can a \emph{fixed}, finite-capacity learner extract from one over
the other? Our bound (\cref{thm:bounded-gap}) is in the same currency these
analyses use --- an approximation term and a Rademacher estimation term
\citep{bartlettmendelson2002, shalevshwartz2014, mohri2018} --- but it carries no
mutual-information term, which is exactly what marks the residual as a capacity
phenomenon rather than an information one.

\subsection{Applied category theory in machine learning}
\label{subsec:rw-cat}

The go/no-go section on parity (\cref{sec:category}) reads operational parity as a
retraction and tracks its round-trip defect. We use category theory only as a
language for that one structural fact, and we want the borrowing to be honest
about its weight. There is a substantial programme that does treat learning
itself categorically --- supervised learning as a functor
\citep{fong2019backpropfunctor} and categorical foundations of gradient-based
learning via lenses and reverse-derivative categories
\citep{cruttwell2022categorical}. We make no contribution to that programme. The
categorical content here is the thin observation that a section/retraction pair
exposes a defect invisible to both the mutual information and the Bayes risk
(\cref{thm:retraction}); the standard references \citep{maclane1998, riehl2016,
awodey2010} supply everything we use, and the section survives only because it
earns a theorem the core sections cannot state (charter §2).

\subsection{Code representation for language models}
\label{subsec:rw-code}

The empirical question this paper formalises is the one studied directly by its
companion preprint, whose related-work section (\texttt{preprint-v1}, §2) we
extend rather than repeat. Two lines bear on it. \emph{Training-time} integration
of code structure bakes ASTs, data flow, or tree structure into model weights ---
CodeBERT \citep{feng2020codebert}, GraphCodeBERT \citep{guo2021graphcodebert},
TreeBERT \citep{jiang2021treebert}, and the generative InCoder
\citep{fried2023incoder}. The closest controlled ancestor is
\citet{namavar2022coderepr}, who compare token, AST, and graph representations for
learning-based program repair and find the better representation is task- and
model-dependent rather than universal. We do not retrain; we vary the input
representation of a frozen model and, unlike that lineage, hold information
constant by construction, so a null is informative rather than inconclusive.

\emph{Inference-time} structure enters at the prompt boundary. Chain-of-thought
\citep{wei2022chainofthought}, plan-then-solve scaffolding
\citep{wang2023plansolve}, and program-aided prompting \citep{gao2023pal} all
improve a fixed model by changing the form, not the information content, of the
input --- and PAL in particular adds an external execution resource, which places
its gains squarely on the capacity/operational side of our account rather than the
information side. A parallel line documents that frontier models are sensitive to
semantically inert prompt-format changes \citep{sclar2024promptformatting,
salinas2024butterfly}; this is precisely the bounded-learner swap gap our positive
half predicts (\cref{sec:bounded}) and our parity-defect section bounds
(\cref{sec:parity-defect}). Finally, two neighbouring areas are sometimes
conflated with representation effects but are operational-consistency mechanisms:
contract- and verifier-backed consistency checking of generated code
\citep{sun2024clover}, and deterministic AST analysis for detecting and correcting
code hallucinations \citep{khati2026asthallucination}. We cite these for what they
are; neither bears on the information-parity claim.

\subsection{Novelty, circumscribed}
\label{subsec:rw-novelty}

Every ingredient above pre-exists: sufficiency and the DPI, the bottleneck
tradeoff, representation-learning bounds, the categorical view of learning, and
the empirical comparison of code formats. The companion preprint's contribution is
a controlled empirical \emph{combination} --- information-parity testing of a
frozen model under pre-registered correction. This paper's contribution is
narrower and theoretical: a formal account of \emph{when} code-representation
format can and cannot help a bounded language model --- never at the information
limit (\cref{thm:irrelevance}), and otherwise only within a capacity budget that
carries no information term and is non-vacuous only in a named finite-resource
regime (\cref{thm:bounded-gap,prop:regime}) --- anchored to a controlled null
rather than to a manufactured positive. We claim the localisation and its
falsifiers (\cref{sec:predictions}), not any of the components it is built from.
